\medskip

\emph{Pour chacune des affirmations suivantes, indiquer si elle est vraie ou fausse.
Chaque réponse doit être justifiée. Une réponse non justifiée ne rapporte aucun point.}

\medskip

\begin{enumerate}
\item Soit $\left(u_n\right)$ une suite définie pour tout entier naturel $n$ et vérifiant la relation suivante:

\begin{center}pour tout entier naturel $n,\:\: \dfrac12  < u_n \leqslant \dfrac{3n^2 + 4n + 7}{6n^2 + 1}$.\end{center}

\textbf{Affirmation 1} $\displaystyle\lim_{n \to + \infty} u_n =\dfrac12$.
\item Soit $h$ une fonction définie et dérivable sur l'intervalle [-4~;~4].

La représentation graphique $\mathcal{C}_{h'}$ de sa fonction dérivée $h'$ est donnée ci-dessous.

\begin{center}
\psset{xunit=0.8cm,yunit=0.4cm,arrowsize=2pt 3}
\begin{pspicture}(-4,-5)(4.5,10)
\psaxes[linewidth=1.25pt,Dy=2](0,0)(-4,-5)(4.5,10)
\psplot[plotpoints=2000,linewidth=1.25pt,linecolor=blue]{-4}{4}{x  4 add  x 3 add mul x 1 add  mul x 3 sub mul  3 x sub mul x 3.85 sub  mul 90 div}
\uput[r](-3.2,2){\blue $\mathcal{C}_{h'}$}
\end{pspicture}
\end{center}

\textbf{Affirmation 2 :} La fonction $h$ est convexe sur $[-1~;~3]$.
\item Le code d'un immeuble est composé de 4 chiffres (qui peuvent être identiques) suivis de deux lettres distinctes parmi A, B et C (exemple: 1232BA).

\textbf{Affirmation 3 :} Il existe 20 634 codes qui contiennent au moins un 0.`

\item On considère la fonction $f$ définie sur $]0~;~+\infty[$ par $f(x) = x \ln x$.

\textbf{Affirmation 4 :}  La fonction $f$ est une solution sur $]0~;~ +\infty[$ de l'équation différentielle

\[xy' - y = x.\]
\end{enumerate}

\bigskip
